\chapter{SOME SPECTRAL THEORY}\label{spectrum}

In this chapter we will, for the first time, assume, unless otherwise stated, that vector
spaces and algebras have complex, rather than real, scalars. A good part of the reason
 \index{conventions!in chapter~\ref{spectrum} scalars are complex}%
for this is the necessity of using \emph{Liouville's theorem}~\ref{C073131} to prove
proposition~\ref{C073134}.

\section{The Spectrum}

\begin{defn} An element $a$ of a unital algebra $A$ is
 \index{invertible!left!in an algebra}%
\df{left invertible} if there exists an element $a_l$ in $A$ (called a \df{left inverse}
of~$a$) such that $a_l a = \vc 1$ and is
 \index{invertible!right!in an algebra}%
\df{right invertible} if there exists an element $a_r$ in $A$ (called a \df{right
inverse} of~$a$)such that $aa_r = \vc 1$. The element is
 \index{invertible!in an algebra}%
\df{invertible} if it is both left invertible and right invertible. The set of all
invertible elements of $A$ is denoted
 \index{invertible@$\inv A$ (invertible elements in an algebra)}%
by~$\inv A$.
\end{defn}

An element of a unital algebra can have at most one multiplicative inverse.  In fact,
more is true.

\begin{prop}\label{000521} If an element of a unital algebra has both a left inverse and a
right inverse, then these two inverses are equal (and so the element is invertible).
\end{prop}

When an element $a$ of a unital algebra is invertible its (unique) inverse is denoted
by~$a^{-1}$.

\begin{prop} If $a$ is an invertible element of a unital algebra, then $a^{-1}$ is also
 \index{inverse!of an inverse}%
invertible and
   \[ \bigl(a^{-1}\bigr)^{-1} = a\,. \]
\end{prop}

\begin{prop} If $a$ and $b$ are invertible elements of a unital algebra, then their product
 \index{inverse!of a product}%
$ab$ is also invertible and
  \[ (ab)^{-1} = b^{-1}a^{-1}\,. \]
\end{prop}

\begin{prop}\label{000524} If $a$ and $b$ are invertible elements of a unital algebra,
 \index{inverses!difference of}%
then
  \[ a^{-1} - b^{-1} = a^{-1}(b - a)b^{-1}\,. \]
\end{prop}

\begin{prop} Let $a$ and $b$ be elements of a unital algebra.  If both $ab$ and
$ba$ are invertible, then so are $a$ and $b$.
\end{prop}

\begin{proof}[\emph{Hint for proof}] Use proposition~\ref{000521}. \ns \end{proof}

\begin{cor} Let $a$ be an element of a unital $*\,$-algebra.  Then $a$ is invertible if and only if $a^*a$ and $aa^*$
are invertible, in which case
   \[ a^{-1} = \bigl(a^*a\bigr)^{-1}a^* = a^*\big(aa^*\bigr)^{-1}\,. \]
\end{cor}

\begin{prop}\label{000526} Let $a$ and $b$ be elements of a unital algebra.  Then
$\vc 1 - ab$ is invertible if and only if $\vc 1 - ba$ is.
\end{prop}

\begin{proof} If $\vc 1 - ab$ is invertible, then $\vc 1 + b(\vc 1 - ab)^{-1}a$ is the
inverse of $\vc 1 - ba$.
\end{proof}

\begin{exam} The simplest example of a complex algebra is the family $\C$ of complex
numbers.  It is both unital and commutative.
\end{exam}

\begin{exam}\label{000532} If $X$ is a nonempty topological space, then the family $\fml C(X)$
of all continuous complex valued functions on $X$ is an
 \index{C@$\fml C(X)$!as a unital algebra}%
algebra under the usual pointwise operations of addition, scalar multiplication, and
multiplication. It is both unital and commutative.
\end{exam}

\begin{exam} If $X$ is a locally compact Hausdorff space which is not compact, then (under
pointwise operations) the family $\fml C_0(X)$ of all continuous complex valued functions
on $X$ which vanish at infinity is a commutative algebra.
However, it is \emph{not} a
 \index{C@$\fml C_0(X)$!as a nonunital algebra}%
unital algebra.
\end{exam}

\begin{exam}\label{000533} The family
 \index{Ma@$\mathbf M_n = \M n\C$!$n \times n$ matrices of complex numbers}%
$\mathbf M_n$ of $n \times n$ matrices of complex numbers is a
 \index{Ma@$\mathbf M_n = \M n\C$!as a unital algebra}%
unital algebra under the usual matrix operations of addition, scalar multiplication, and
multiplication. It is not commutative when $n > 1$.
\end{exam}

\begin{exam}\label{000533a} Let $A$ be an algebra. Make the family
 \index{Ma@$\M nA$!$n \times n$ matrices of elements of an algebra}%
$\M n{A}$ of $n \times n$ matrices of elements of $A$ into an
 \index{Ma@$\M nA$!as an algebra}%
algebra by using the same rules for matrix operations that are used for~$\mathbf M_n$.
Thus $\mathbf M_n$ is just $\M n\C$. The algebra $\M nA$ is unital if and only if $A$~is.
\end{exam}

\begin{exam}\label{000534} If $V$ is a normed linear space, then
 \index{B@$\ofml B(V)$!as a unital algebra}%
$\ofml B(V)$ is a unital algebra.  If $\dim V~>~1$, the algebra is not commutative.
\end{exam}

\begin{defn}\label{000535} Let $a$ be an element of a unital algebra~$A$.  The
 \index{spectrum}%
\df{spectrum} of $a$, denoted by
 \index{sigma@$\sigma_A(a)$ or $\sigma(a)$ (spectrum of~$a$)}%
$\sigma_A(a)$ or just $\sigma(a)$, is the set of all complex numbers $\lambda$ such that
$a - \lambda\vc 1$ is not invertible.
\end{defn}

\begin{exam} If $z$ is an element of the algebra $\C$ of complex numbers, then
 \index{spectrum!of a complex number}%
$\sigma(z) = \{z\}$.
\end{exam}

\begin{exam}\label{0005712} Let $X$ be a compact Hausdorff space.  If $f$ is an element
 \index{spectrum!of a continuous function}%
of the algebra $\fml C(X)$ of continuous complex valued functions on $X$, then the
spectrum of $f$ is its range.
\end{exam}

\begin{exam} Let $X$ be an arbitrary topological space.  If $f$ is an element of the
 \index{spectrum!of a bounded continuous function}%
algebra $\fml C_b(X)$ of bounded continuous complex valued functions on $X$, then the
spectrum of $f$ is the closure of its range.
\end{exam}

\begin{exam} Let $S$ be a positive measure space and $f \in \lfs{\infty}{S}$ be an
(equivalence class of) essentially bounded function(s) on~$S$. Then the
 \index{spectrum!of an essentially bounded function}%
spectrum of $f$ is its essential range.
\end{exam}

\begin{exam} The family $\mathbf M_3$ of $3 \times 3$ matrices of complex numbers is a unital
algebra under the usual matrix operations. The spectrum of the matrix $\begin{bmatrix} 5
& -6 & -6 \\ -1 & 4 & 2 \\ 3 & -6 & -4 \end{bmatrix}$ is $\{1,2\}$.
\end{exam}

\begin{prop} Let $a$ be an element of a unital algebra such that $a^2 = \vc 1$.
 \index{spectrum!of an element whose square is $1$}%
Then either
 \begin{enumerate}[label=\rm{(\alph*)}]
  \item $a = \vc 1$, in which case $\sigma(a) = \{1\}$, or
  \item $a = -\vc 1$, in which case $\sigma(a) = \{-1\}$, or
  \item $\sigma(a) = \{-1,1\}$.
 \end{enumerate}
\end{prop}

\begin{proof}[\emph{Hint for proof}]  In (c) to prove $\sigma(a) \subseteq \{-1,1\}$,
consider $\D\frac1{1 - \lambda^2}(a + \lambda \mathbf 1).$  \ns
\end{proof}

\begin{prop} Recall that an element $a$ of an algebra is
 \index{idempotent}%
 \index{spectrum!of an idempotent element}%
\df{idempotent} if $a^2 = a$. Let $a$ be an idempotent element of a unital algebra. Then
either
 \begin{enumerate}[label=\rm{(\alph*)}]
  \item $a = \vc 1$, in which case $\sigma(a) = \{1\}$, or
  \item $a = \vc 0$, in which case $\sigma(a) = \{0\}$, or
  \item $\sigma(a) = \{0,1\}$.
 \end{enumerate}
\end{prop}

\begin{proof}[\emph{Hint for proof}] In (c) to prove $\sigma(a) \subseteq \{0,1\}$,
consider $\D \frac1{\lambda - \lambda^2}\bigl(a + (\lambda - 1)\vc 1\bigr).$ \ns
\end{proof}

\begin{prop} Let $a$ be an invertible element of a unital algebra. Then a complex number
$\lambda$ belongs to the spectrum of $a$ if and only if $1/\lambda$ belongs to the
spectrum of~$a^{-1}$.
\end{prop}

The next proposition is a simple corollary of proposition~\ref{000526}.

\begin{prop}\label{000575} If $a$ and $b$ are elements of a unital algebra, then,
except possibly for $0$, the spectra of $ab$ and $ba$ are the same.
\end{prop}

\begin{defn}\label{C035854} A map $f \colon A \sto B$ between Banach algebras is a
 \index{Banach!algebra!homomorphism}%
 \index{homomorphism!of Banach algebras}%
\df{(Banach algebra) homomorphism} if it is both an algebraic homomorphism and a
continuous map between $A$ and $B$.  We denote
 \index{BALG@$\cat{BALG}$!the category}%
 \index{category!$\cat{BALG}$ as a}%
by $\cat{BALG}$ the category of Banach algebras and continuous algebra homomorphisms.
\end{defn}

\begin{prop} Let $A$ and $B$ be Banach algebras.  The
 \index{direct!sum!of Banach algebras}%
 \index{Banach!algebra!direct sum}%
\df{direct sum} of $A$ and $B$, denoted by $A \oplus B$, is defined to be the set $A
\times B$ equipped with pointwise operations:
   \begin{enumerate}[label=\rm{(\alph*)}]
     \item $(a,b) + (a',b') = (a + a',b + b')$;
     \item $(a,b)\,(a',b') = (aa',bb')$;
     \item $\alpha(a,b) = (\alpha a,\alpha b)$
   \end{enumerate}
for $a$,$a' \in A$, $b$, $b' \in B$, and $\alpha \in \C$. As in~\ref{C023414} define
$\norm{(a,b)} = \norm a + \norm b$.  This makes $A \oplus B$ into a Banach algebra.
\end{prop}

\begin{exer} Identify products and coproducts (if they exist) in the category $\cat{BALG}$
of Banach algebras
 \index{product!in $\cat{BALG}$}%
 \index{Banach!algebra!product}%
 \index{Banach!algebra!coproduct}%
 \index{coproduct!in $\cat{BALG}$}%
 \index{BALG@$\cat{BALG}$!products and coproducts in}%
and continuous algebra homomorphisms and in the category of Banach algebras and
contractive algebra homomorphisms. (Here
 \index{contractive}%
\emph{contractive} means $\norm{Tx} \le \norm x$ for all~$x$.)
\end{exer}

\begin{prop} In a Banach algebra $A$ the operations of addition and multiplication (regarded as
maps from $A \oplus A$ to $A$) are continuous and scalar multiplication (regarded as a map from
$\K \oplus A$ to $A$) is also continuous.
\end{prop}

\begin{prop}[The Neumann series]\label{prop_Neumann_series} Let $a$ be an element of a unital Banach algebra.
 \index{Neumann series}%
 \index{series!Neumann}%
If $\norm a < 1$, then
$\vc 1 - a \in \inv A$ and $(\vc 1 - a)^{-1} = \sum_{k=0}^\infty a^k$.
\end{prop}

\begin{proof}[\emph{Hint for proof}]  Show that the sequence $(\vc 1, a, a^2, \dots)$ is
summable.  (In a unital algebra we take $a^0$ to mean~$\vc 1$.)  \ns
\end{proof}

\begin{cor}\label{cor_Neumann_series} If $a$ is an element of a unital Banach algebra with $\norm{\vc 1 - a} < 1$, then $a$ is invertible and
    \[ \norm{a^{-1}} \le \frac{1}{1 - \norm{\vc 1 - a}}\,.  \]
\end{cor}

\begin{cor}\label{cor2_Neumann_series} Suppose that $a$ is an invertible element of a unital Banach algebra $A$, that $b \in A$,
and that $\norm{a - b} < \norm{a^{-1}}^{-1}$.  Then $b$ is invertible in~$A$.
\end{cor}

\begin{proof}[\emph{Hint for proof}] Use the preceding corollary to show that $a^{-1}b$ is invertible.
\ns \end{proof}

\begin{cor} If $a$ belongs to a unital Banach algebra and $\norm a < 1$, then
   \[ \norm{(\vc 1 - a)^{-1} - \vc 1} \le \frac{\norm a}{1 - \norm a}\,. \]
\end{cor}

\begin{prop}\label{000644fa} If $A$ is a unital Banach algebra, then $\open{\inv A}A$.
\end{prop}

\begin{proof}[\emph{Hint for proof}] Let $a \in \inv A$.  Show, for sufficiently small $h$,
that $\vc 1 - a^{-1}h$ is invertible.  \ns
\end{proof}

\begin{prop}\label{000646fa} Let $A$ be a unital Banach algebra.  The map $a \mapsto a^{-1}$ from $\inv A$ into
itself is a homeomorphism.
\end{prop}

\begin{prop} Let $A$ be a unital Banach algebra.  The map $r\colon a \mapsto a^{-1}$ from
$\inv A$ into itself is differentiable and at each invertible element $a$, we have
$dr_a(h) = -a^{-1}ha^{-1}$ for all $h \in A$.
\end{prop}

\begin{proof}[\emph{Hint for proof}] For a quick introduction to differentiation on infinite dimensional
spaces, see Chapter 13 of my notes~\cite{Erdman:2007} on Real Analysis.  \ns
\end{proof}

\begin{prop}\label{000661} Let $a$ be an element of a unital Banach algebra~$A$.  Then the spectrum of $a$
is compact and $\abs\lambda \le \norm a$ for every $\lambda \in \sigma(a)$.
\end{prop}

\begin{proof}[\emph{Hint for proof}] Use the \emph{Heine-Borel theorem}.  To prove that the
spectrum is closed notice that $(\sigma(a))^c = f^\gets(\inv A)$ where $f(\lambda) = a -
\lambda\mathbf 1$ for every complex number~$\lambda$. Also show that if $\abs\lambda >
\norm a$, then $\mathbf 1 - \lambda^{-1}a$ is invertible.  \ns
\end{proof}

\begin{defn}  Let $a$ be an element of a unital Banach algebra.  The
 \index{resolvent!mapping}%
\df{resolvent mapping} for $a$ is defined by
   \[ R_a \colon \C \setminus \sigma(a) \sto A
                 \colon \lambda \mapsto (a  - \lambda\vc 1)^{-1}\,. \]
\end{defn}

\begin{defn} Let $\open U\C$ and $A$ be a unital Banach algebra. A function $f \colon U \sto A$
 \index{analytic}%
is \df{analytic} on $U$ if
   \[ f'(a) := \lim_{z \sto a} \frac{f(z) - f(a)}{z - a} \]
exists for every $a \in U$.  A complex valued function which is analytic on all of $\C$
is an
 \index{entire}
\df{entire} function.
\end{defn}

\begin{prop} For $a$ an element of a unital Banach algebra $A$ and $\phi$ a bounded linear
functional on~$A$ let $f := \phi \circ R_a \colon \C \setminus \sigma(a) \sto \C$.  Then
 \begin{enumerate}[label=\rm{(\alph*)}]
  \item $f$ is analytic on its domain, and
  \item $f(\lambda) \sto 0$ as $\abs\lambda \sto \infty$.
 \end{enumerate}
\end{prop}

\begin{proof}[\emph{Hint for proof}] For (i) notice that in a unital algebra $a^{-1} - b^{-1}
= a^{-1}(b - a)b^{-1}$.   \ns
\end{proof}

\begin{thm}[Liouville]\label{C073131} Every bounded entire function on $\C$ is constant.
 \index{Liouville's theorem}%
\end{thm}

\begin{proof}[\emph{Comment on proof}] The proof of this theorem requires a little background
in the theory of analytic functions, material that is not covered in these notes.  The
theorem would surely be covered in any first course in complex variables.  For example, you can find a
nice proof in \cite{Rudin:1987}, theorem 10.23.   \ns
\end{proof}

\begin{prop}\label{C073134} The spectrum of every element of a unital Banach algebra is nonempty.
\end{prop}

\begin{proof}[\emph{Hint for proof}] Argue by contradiction.  Use \emph{Liouville's theorem}~\ref{C073131}
to show that $\phi \circ R_a$ is constant for every bounded linear functional $\phi$ on~$A$. Then use
(a version of) the \emph{Hahn-Banach theorem} to prove that $R_a$ is constant.  Why must this constant be~$0$?   \ns
\end{proof}

\noindent\emph{Note.} It is important to keep in mind that we are working only with
\emph{complex} algebras.  This result is \emph{false} for real Banach algebras.  An easy
counterexample is the matrix $\begin{bmatrix}0 & 1 \\ -1 & 0 \end{bmatrix}$ regarded as
an element of the (real) Banach algebra $M_2$ of all $2 \times 2$ matrices of real
numbers.

 \vskip 1 pt

\begin{thm}[Gelfand-Mazur]\label{C073141} If $A$ is a unital Banach algebra in which every
 \index{Gelfand-Mazur theorem}%
nonzero element is invertible, then $A$ is isometrically isomorphic to~$\C$.
\end{thm}

\begin{proof}[\emph{Hint for proof}] Let $B = \{ \lambda\vc 1\colon \lambda \in \C\}$. Use
the preceding proposition~\ref{C073134}) to show that $B = A$.    \ns
\end{proof}

The following is an immediate consequence of the \emph{Gelfand-Mazur theorem}
\ref{C073141} and proposition~\ref{C037821}.

\begin{cor}\label{C073147} If $I$ is a maximal ideal in a commutative unital Banach algebra $A$,
then $A/I$ is isometrically isomorphic to~$\C$.
\end{cor}

\begin{prop}\label{0006703fa} Let $a$ be an element of a unital algebra.  Then
   \[ \sigma(a^n) = [\sigma(a)]^n \]
for every $n \in \N$.  (The notation $[\sigma(a)]^n$ means $\{\lambda^n\colon \lambda \in
\sigma(a)\}$.)
\end{prop}

\begin{defn}\label{C073154}  Let $a$ be an element of a unital algebra.  The
 \index{spectral!radius}%
 \index{radius!spectral}%
\df{spectral radius} of~$a$, denoted by $\rho(a)$, is defined to be
$\sup\{\abs\lambda\colon \lambda \in \sigma(a)\}$.
\end{defn}

\begin{prop}\label{0006703} If $a$ is an element of a unital Banach algebra, then $\rho(a) \le
\norm a$ and $\rho(a^n) = \bigr(\rho(a)\bigr)^n$.
\end{prop}

\begin{prop} Let $\phi\colon A \sto B$ be a unital homomorphism between unital algebras. If $a$ is an invertible element of $A$,
then $\phi(a)$ is invertible in $B$ and its inverse is~$\phi(a^{-1})$.
\end{prop}

\begin{cor} If $\phi\colon A \sto B$ is a unital homomorphism between unital algebras, then
   \[ \sigma(\phi(a)) \subseteq \sigma(a) \]
for every $a \in A$.
\end{cor}

\begin{cor} If  $\phi\colon A \sto B$ is a unital homomorphism between unital algebras, then
   \[ \rho(\phi(a)) \le \rho(a) \]
for every $a \in A$.
\end{cor}

\begin{thm}[Spectral radius formula]\label{C073157} If $a$ is an element of a unital Banach
 \index{spectral!radius!formula for}%
algebra, then
   \[ \rho(a) = \inf\bigl\{\norm{a^n}^{1/n}\colon n \in \N\bigr\}
                                = \lim_{n \sto \infty} \norm{a^n}^{1/n}\,. \]
\end{thm}

\begin{exer} The Volterra operator $V$ (defined in example~\ref{000319}) is an
 \index{Volterra operator!and integral equations}%
 \index{operator!Volterra!and integral equations}%
integral operator on the Banach space~$\fml C([0,1])$.
 \begin{enumerate}
  \item[(a)] Compute the spectrum of the operator $V$.  (\emph{Hint.} Show that $\norm{V^n}
\le (n!)^{-1}$ for every $n \in \N$ and use the \emph{spectral radius formula}.)
  \item[(b)] What does (a) say about the possibility of finding a continuous function $f$ which
satisfies the integral equation
   \[ f(x) - \mu\int_0^xf(t)\,dt = h(x), \]
where $\mu$ is a scalar and $h$ is a given continuous function on~$[0,1]$?
  \item[(c)] Use the idea in (b) and the Neumann series expansion for $\vc 1 - \mu V$ (see
proposition~\ref{prop_Neumann_series}) to calculate explicitly (and in terms of elementary functions)
a solution to the integral equation
 \[f(x) - {\textstyle \frac12}\int_0^xf(t)\,dt = e^x.\]
 \end{enumerate}
\end{exer}

\begin{prop} Let $A$ be a unital Banach algebra and $B$ a closed subalgebra containing~$\vc 1_A$.
Then
 \begin{enumerate}[label=\rm{(\alph*)}]
   \item $\inv B$ is both open and closed in $B \cap \inv A$;
   \item $\sigma_A(b) \subseteq \sigma_B(b)$ for every $b \in B$; and
   \item if $b \in B$ and $\sigma_A(b)$ has no holes (that is, if its complement in $\C$ is
connected), then $\sigma_A(b) =    \sigma_B(b)$.
 \end{enumerate}
\end{prop}

\begin{proof}[\emph{Hint for proof}] Consider the function
$f_b\colon (\sigma_A(b))^c \sto B \cap \inv A\colon \lambda \mapsto b - \lambda \vc 1$.   \ns
\end{proof}















\section{Spectra of Hilbert Space Operators}
\begin{prop} For a self-adjoint operator $T$ on a Hilbert space the following are
equivalent:
 \begin{enumerate}
    \item[(i)] $\sigma(T) \subseteq [0,\infty)$;
    \item[(ii)] there exists an operator $S$ on $H$ such that $T = S^*S$; and
    \item[(iii)] $T$ is positive.
 \end{enumerate}
\end{prop}

\begin{notn} If $A$ is a set of the complex numbers we
 \index{<unopurstar@$A^*$ (set of complex conjugates of $A \subseteq \C$)}%
define $A^* := \{\conj\lambda \colon \lambda \in A\}$.  We do this to avoid confusion with the closure of the set~$A$.
\end{notn}

\begin{prop} If $T$ is a Banach space operator then $\sigma(T^*) = \sigma(T)$, while if it is a Hilbert space
operator then $\sigma(T^*) = (\sigma(T))^*$.
\end{prop}

In a finite dimensional space there is only one way that a complex number $\lambda$ can end up in the spectrum of an operator~$T$.
The reason for this is that there is only one way an operator, in this case $T - \lambda I$, on a finite dimensional space can fail
to be invertible (see the paragraph before definition~\ref{defn_similar_vs_ops}).  On the other hand there are several ways for $\lambda$
to get into the spectrum of an operator on an infinite dimensional space---because an operator has several ways in which it may fail
to be invertible. (See, for example, proposition~\ref{prop_nasc_op_invert}.) Historically this has led to a number of proposed
taxonomies of the spectrum of an operator.

\begin{defn} Let $T$ be an operator on a Hilbert (or Banach) space~$H$. A complex number $\lambda$ belongs to the
 \index{resolvent!set}%
\df{resolvent set} of $T$ if $T - \lambda I$ is invertible. The resolvent set is the complement of the \emph{spectrum}
of $T$. The set of complex numbers $\lambda$ for which $T - \lambda I$ fails to be injective is the
 \index{point spectrum}%
 \index{spectrum!point}%
\df{point spectrum} of $T$ and is denoted
 \index{sigma@$\sigma_p(T)$ (point spectrum of~$T$)}%
by $\sigma_p(T)$. Members of the point spectrum of $T$ are
 \index{eigenvalue}%
\df{eigenvalues} of $T$. The set of complex numbers $\lambda$ for which $T - \lambda I$ is not bounded away from zero is the
 \index{approximate point spectrum}%
 \index{spectrum!approximate point}%
\df{approximate point spectrum} of $T$ and is denoted
 \index{sigma@$\sigma_{ap}(T)$ (approximate point spectrum of~$T$)}%
by $\sigma_{ap}(T)$. The set of all complex numbers $\lambda$ such that the closure of the range of $T - \lambda I$ is a proper
subspace of $H$ is the
 \index{compression spectrum}%
 \index{spectrum!compression}%
\df{compression spectrum} of $T$ and is denoted
 \index{sigma@$\sigma_c(T)$ (compression spectrum of~$T$)}%
by $\sigma_c(T)$.  The
 \index{residual spectrum}%
 \index{spectrum!residual}%
\df{residual spectrum} of $T$, which we denote
 \index{sigma@$\sigma_r(T)$ (residual spectrum of~$T$)}%
by $\sigma_r(T)$, is $\sigma_c(T) \setminus \sigma_p(T)$.
\end{defn}

\begin{prop} If $T$ is a Hilbert space operator, then $\sigma_p(T) \subseteq \sigma_{ap}(T) \subseteq \sigma(T)$.
\end{prop}

\begin{prop} If $T$ is a Hilbert space operator, then $\sigma(T) = \sigma_{ap}(T) \cup \sigma_c(T)$.
\end{prop}

\begin{prop} If $T$ is a Hilbert space operator, then $\sigma_p(T^*) = (\sigma_c(T))^*$.
\end{prop}

\begin{cor} If $T$ is a Hilbert space operator, then $\sigma(T) = \sigma_{ap}(T) \cup (\sigma_{ap}(T^*))^*$.
\end{cor}

\begin{thm}{Spectral Mapping Theorem}\label{SMThm} If $T$ is a Hilbert space operator and $p$ is a polynomial, then
   \[ \sigma(p(T)) = p(\sigma(T)). \]
\end{thm}

The preceding theorem is true more generally for any analytic function defined in a neighborhood of the spectrum of~$T$.
(For a proof see~\cite{Conway:1990}, chapter VII, theorem 4.10.)  The result also holds for parts of the spectrum.

\begin{prop} If $T$ is a Hilbert space operator and $p$ is a polynomial, then
   \begin{enumerate}[label=\rm{(\alph*)}]
     \item $\sigma_p(p(T)) = p(\sigma_p(T))$,
     \item $\sigma_{ap}(p(T)) = p(\sigma_{ap}(T))$, and
     \item $\sigma_c(p(T)) = p(\sigma_c(T))$,
   \end{enumerate}
\end{prop}

\begin{prop}
    If $T$ is an invertible operator on a Hilbert space, then  $\sigma(T^{-1}) = (\sigma(T))^{-1}$.
\end{prop}

\begin{defn}\label{defn_similar_Hsp_ops} As with vector spaces, two operators $R$ and $T$ on a Hilbert (or Banach) space $H$ are
 \index{similar!operators}%
\df{similar} if there exists an invertible operator $S$ on $H$ such that $R = STS^{-1}$.
\end{defn}

\begin{prop} If $S$ and $T$ are similar operators on a Hilbert space, then $\sigma(S) = \sigma(T)$, $\sigma_p(S) = \sigma_p(T)$,
$\sigma_{ap}(S) = \sigma_{ap}(T)$, and $\sigma_c(S) = \sigma_c(T)$.
\end{prop}

\begin{prop} If $T$ is a normal operator on a Hilbert space, then $\sigma_c(T) = \sigma_p(T)$, so that $\sigma_r(T)~=~\emptyset$.
\end{prop}

\begin{prop} If $T$ is an operator on a Hilbert space, then $\sigma_{ap}(T)$ is closed.
\end{prop}

\begin{exam} Let $D = \diag(a_1,a_2,a_3,\dots)$ be a diagonal operator on a Hilbert space and let
 \index{spectral!parts!of a diagonal operator}%
 \index{operator!diagonal!spectral parts of}%
 \index{diagonal!operator!spectral parts of}%
$A = \cup_{k=1}^\infty a_k$. Then
   \begin{enumerate}[label=(\alph*)]
     \item $\sigma_p(D) = \sigma_c(D) = A$,
     \item $\sigma_{ap}(D) = \clo A$, and
     \item $\sigma_r(D) = \emptyset$.
   \end{enumerate}
\end{exam}

\begin{exam} Let $S$ be the unilateral shift operator on $l_2$
 \index{spectral!parts!of the unilateral shift}%
 \index{operator!unilateral shift!spectral parts of}%
 \index{unilateral shift!spectral parts of}%
(see example~\ref{C063526}). Then
    \begin{enumerate}[label=(\alph*)]
      \item $\sigma(S) = D$,
      \item $\sigma_p(S) = \emptyset$,
      \item $\sigma_{ap}(S) = C$,
      \item $\sigma_c(S) = B$,
      \item $\sigma(S^*) = D$,
      \item $\sigma_p(S^*) = B$,
      \item $\sigma_{ap}(S^*) = D$, and
      \item $\sigma_c(S^*) = \emptyset$.
    \end{enumerate}
where $B$ is the open unit disk, $D$ is the closed unit disk, and $C$ is the unit circle in the complex plane.
\end{exam}

\begin{exam} Let $(S,\sfml A, \mu)$ be a $\sigma$-finite positive measure space, $\phi \in L_\infty(S,\mu)$, and $M_\phi$ be the
 \index{spectral!parts!of a multiplication operator}%
 \index{operator!multiplication!spectral parts of}%
 \index{multiplication operator!spectral parts of}%
corresponding multiplication operator on the Hilbert space $\lfs2{S,\mu}$ (see example~\ref{C063527}). Then
the spectrum and the approximate point spectrum of $M_\phi$ are the essential range of~$\phi$, while the point spectrum
and compression spectrum of $M_\phi$ are $\{\lambda \in \C\colon \mu(\phi^{\gets}(\{\lambda\})) > 0\}$.
\end{exam}






















\endinput








